\documentclass[11pt]{amsart}
\usepackage[T1]{fontenc}
\usepackage[margin=1.02in]{geometry}
\DeclareFontShape{OT1}{cmr}{bx}{n}{<-> cmbx10}{}
\DeclareMathSizes{8}{8}{8}{8}
\DeclareMathSizes{9}{9}{8}{8}
\DeclareMathSizes{10.95}{10.95}{8}{8}
\newtheorem{theorem}{Theorem}
\title{An endpoint converse for order-zero Bessel--Fourier coefficients}
\date{November 27, 2026}
\begin{document}
\begin{abstract}
Given any nonnegative null sequence \(a_m\), a nonnegative
function \(f\) in \(L^1([0,1],x\,dx)\), supported in
\([0,1/2]\), has order-zero Bessel--Fourier coefficients
with \(a_m\le f_m/\sqrt{j_m}\to0\). The construction uses
dilates of a radial autocorrelation and dyadically chosen
masses. It proves Acosta Babb's endpoint existence statement
unconditionally. A separate heat-semigroup argument rules
out the stronger single-parity statement.
\end{abstract}
\maketitle

\section{The endpoint and its kernel}
Let \(j_m\) be the \(m\)-th positive zero of \(J_0\), with
the coefficient convention of Acosta Babb \cite{acosta},
\[
 f_m=\frac{2}{J_1(j_m)^2}\int_0^1 f(x)J_0(j_m x)x\,dx .
\]
The endpoint existence statement is Conjecture~3 of the
dissertation \cite[Sec.~5.5, p.~88]{thesis} and appears as
a conditional Corollary~5.2 in the journal version
\cite[p.~13]{acosta}.

\begin{theorem}
For every sequence \(a_m\ge0\) with \(a_m\to0\), there is a
nonnegative \(f\in L^1([0,1],x\,dx)\), supported in
\([0,1/2]\), such that
\[
                  0\le a_m\le f_m/\sqrt{j_m}\to0 .
\]
\end{theorem}

Choose a radial nonnegative smooth probability density \(h\)
on \(\mathbb R^2\) supported in the disk of radius \(1/2\).
Set \(g=h*\widetilde h\), with \(\widetilde h(z)=h(-z)\).
The function \(g\) is radial, nonnegative, smooth, of integral
one, and supported in the unit disk. Using
\(\widehat g(\xi)=\int_{\mathbb R^2}e^{-i\xi\cdot z}g(z)\,dz\),
write \(\widehat g(\xi)=\psi(|\xi|)\). Then
\(\psi=|\widehat h|^2\ge0\), \(\psi(0)=1\), and \(\psi\le1\).
Integration by parts gives rapid decay. Consequently there
are \(s_0\in(0,1/2]\) and \(C>0\) such that
\[
  \psi(s)\ge\frac12\quad (0\le s\le s_0),\qquad
  \psi(s)\le C\min(1,s^{-2})\quad(s>0).
\]
For \(0<r\le1\) set \(g_r(z)=r^{-2}g(z/r)\) and
\(F_r(x)=2\pi g_r((x,0))\). Polar coordinates give
\[
 \int_0^1 F_r(x)x\,dx=1,\qquad
 \int_0^1 F_r(x)J_0(j_m x)x\,dx=\psi(r j_m).
\]
For the second identity integrate the Fourier phase in the
angular variable to obtain \(J_0\). The same angular
representation gives \(|J_0(y)|\le1\) for real \(y\).

\section{Choice of scales and both estimates}
The Bessel asymptotics in \cite[Lemma~2.2]{acosta} give
\(j_m\asymp m\) and \(|J_1(j_m)|\asymp j_m^{-1/2}\).
The finitely many small indices are included since the positive
zeros of \(J_0\) are simple and \(J_0'=-J_1\).
Thus
\[
   B_m=\frac{2}{\sqrt{j_m}J_1(j_m)^2}\asymp\sqrt m.
\]
Fix \(c>0\), \(D\ge1\) with \(B_m\ge c\sqrt m\) and
\(j_m\le Dm\) for all \(m\). For \(k\ge0\) set
\[
  A_k=\max_{2^k\le m<2^{k+1}}a_m,\qquad
  r_k=\frac{s_0}{D\,2^{k+1}},\qquad
  w_k=\frac2c\,A_k2^{-k/2}.
\]
The sequence \(A_k\) is bounded and tends to zero, so
\(\sum_k w_k<\infty\). The nonnegative sum
\(f=\sum_{k\ge0}w_kF_{r_k}\) converges in the weighted
\(L^1\) space, with support in \([0,1/2]\).
Since \(J_0\) is bounded, each coefficient may be taken
term by term. In particular
\[
                  f_m/\sqrt{j_m}=B_m\sum_{k\ge0}w_k\psi(r_k j_m).
\]
Each summand is nonnegative. If \(m\) is in block \(\ell\),
then \(r_\ell j_m\le s_0\) and
\(B_m\ge c\,2^{\ell/2}\). The summand \(k=\ell\) alone
contributes at least \(A_\ell\ge a_m\).

For convergence, \(B_m\le C_1 2^{\ell/2}\), and
\(r_kj_m\ge C_2 2^{\ell-k}\) for absolute constants determined
by \(h\) and the preceding comparisons. For \(k\ge\ell\) use
\(\psi\le1\). For \(k<\ell\) use the inverse-square bound.
The normalized coefficient is at most a fixed multiple of
\[
       \sum_{k\ge\ell} A_k2^{-(k-\ell)/2}
         + \sum_{k<\ell} A_k2^{-3(\ell-k)/2}.
\]
The first sum is bounded by a geometric constant times
\(\sup_{k\ge\ell}A_k\). In the second, each fixed initial term
decays as \(\ell\to\infty\), and the tail is bounded by a
geometric constant times \(\sup_{k\ge K} A_k\) for any cutoff
\(K\). Both sums tend to zero. This proves the theorem.

\section{The single-parity obstruction}
The same article proposes in Conjecture~2 that one parity may
vanish while the other has the stated positive lower bounds
\cite[p.~13]{acosta}. This stronger statement fails for
\(a_m=1/\log(m+2)\), even without a decay requirement.
Both index conventions in its two clauses give the same
growth estimate on the retained parity.

Let \(U_t\) be radial Dirichlet heat evolution on the unit disk
for \(t>0\). The heat semigroup preserves radial functions.
It is positive and \(L^1\)-contractive: its symmetric kernel is
nonnegative and has mass at most one in each variable, by the
killed heat representation or the maximum principle and
symmetry \cite{heat}. In radial coordinates,
\[
 \left|\int_0^1 (U_t f)(x)x\,dx\right|
     \le\|f\|_{L^1(x\,dx)} .
\]
The radial eigenfunctions are \(J_0(j_m x)\) with heat factors
\(e^{-tj_m^2}\), and their integral against \(x\,dx\) is
\(J_1(j_m)/j_m\) \cite[Eq.~(2.8)]{acosta}. Therefore
\[
       \int_0^1 (U_t f)(x)x\,dx
         =\sum_{m\ge1} e^{-tj_m^2} f_m\frac{J_1(j_m)}{j_m}.
\]
For bounded radial \(f\) this is the \(L^2\) expansion paired
with the constant function. For arbitrary \(f\) in the weighted
\(L^1\) space use bounded truncations. Contractivity gives
convergence on the left, and on the right the uniform bound
\(|f_m|\le C_3 j_m\|f\|_{L^1(x\,dx)}\) gives absolute
summability for every \(t>0\) and justifies passage to the
limit.

The signs of \(J_1(j_m)\) alternate, since each zero of \(J_0\)
is simple and \(J_0'=-J_1\). Thus the eventual nonzero terms
have the same sign. The positive lower bounds on the retained
parity imply their unsigned magnitudes without the heat factor
are at least a constant times \(1/(m\log(m+2))\). The latter
sum diverges on either parity. After discarding finitely many
terms, whose heat contributions stay bounded, monotone
convergence as \(t\downarrow0\) contradicts the heat contraction.
The source's conditional Proposition~5.3 already notes the
same-sign unweighted integral of the partial sums. The
heat argument turns that behavior into an obstruction to the
hypothesized coefficient pattern. No assertion about general
unregularized divergence follows.

\section*{Prior results and verification}
Acosta Babb's Theorem~1.4 covers exponents below \(1/2\);
Corollary~5.2 uses the separate conjecture about positivity of
a modified Fej\'er kernel. Reports by TruReTuring and Kang
\cite{tru,kang} give two-term counterexamples to that positivity
assertion. They leave the endpoint existence statement
undecided. Here the autocorrelation multiplier is nonnegative
by construction. The dyadic weighting then compensates for
the Bessel factor \(\sqrt m\).

The replay checks the exact dyadic algebra and finite numerical
controls for the illustration. A partial Lean~4 certificate
checks the two squared-envelope power identities, an own-block
index comparison, and an exact scalar floor used by the
illustration. The analytic existence and heat obstruction
are written proofs. Internally reviewed, not peer reviewed.

Under human direction, ChatGPT Astra~6.2 (OpenAI) generated the
discovery proof and research packet. OpenAI Codex independently
checked the mathematics. GPT-6.2 Sol (OpenAI) checked the
literature and prepared the manuscript, exact dyadic replay,
partial Lean certificate and illustrations. DannyExperiments
maintains the repository and is the preferred citation identity.
No AI system is proposed as an author.

\begin{thebibliography}{9}
\raggedright
\bibitem{thesis} R. L. Acosta Babb, \emph{\(L^p\) convergence of
eigenfunction expansions for the Laplacian on planar domains},
PhD thesis, University of Warwick, 2024.
\bibitem{acosta} R. L. Acosta Babb,
\emph{On the rate of growth of Bessel--Fourier coefficients for integrable
functions}, J. Math. Anal. Appl. 559(1) (2026), 130440.
doi:10.1016/j.jmaa.2026.130440.
\bibitem{heat} E. B. Davies and B. Simon,
\emph{Ultracontractivity and the heat kernel for Schr\"odinger operators and
Dirichlet Laplacians}, J. Funct. Anal. 59 (1984), 335--395.
\bibitem{tru} TruReTuring, \emph{Acosta Babb's weighted Bessel-kernel
positivity conjecture}, public research dossier ab2410.
\bibitem{kang} M. Kang, \emph{A two-term counterexample to Bessel--Fourier
kernel positivity}, MathDB 2610.00793v1, 2026.
\end{thebibliography}
\end{document}
